% Guia do professor — Capítulo 17: Ventos Regionais
\section{Capítulo 17 --- Ventos Regionais}

\subsection*{Visão geral e ritmo}

O capítulo da geografia local: brisas, vales, catabáticos, ondas de
montanha e a hidráulica de sotavento. Fisicamente delicioso — teorema
de circulação, oscilador $N$, água rasa e Froude — e imediatamente
reconhecível pela turma (Santos, serra do Mar, vale do Paraíba).
\textbf{Ritmo: 2 aulas de 2\,h.}

\begin{itemize}
  \item \textbf{Aula 1} — Bjerknes e a brisa (com o hodógrafo do
        Caso~1), terral, ventos de vale/encosta, catabáticos até a
        Antártida (Caso~2).
  \item \textbf{Aula 2} — ondas de montanha e Scorer (Caso~3), föhn
        completo, água rasa/Froude e vendavais (Caso~4), gap winds.
\end{itemize}

\subsection*{Objetivos de aprendizagem}

\begin{enumerate}
  \item explicar a brisa pelo teorema de circulação e prever seu giro
        diurno no HS;
  \item calcular a velocidade de equilíbrio catabática;
  \item usar o parâmetro de Scorer para classificar ondas
        (propagante/evanescente/aprisionada) e obter $\lambda_z$;
  \item calcular o aquecimento föhn com remoção de chuva;
  \item usar o número de Froude para prever bloqueio, travessia e
        vendaval de sotavento.
\end{enumerate}

\subsection*{Pré-requisitos a reativar}

Bowen e contraste terra--mar (Cap.~3); resfriamento radiativo noturno
(Caps.~2, 6); $N$ e estabilidade (Cap.~5); $\theta$ e föhn (Cap.~3);
nuvens de onda (Cap.~6); Coriolis (Cap.~10).

\subsection*{Armadilhas conceituais frequentes}

\begin{itemize}
  \item \textbf{``Brisa é o ar indo do frio pro quente''}: embaixo
        sim, mas é um CIRCUITO — sem o ramo de retorno em altitude
        não fecha; o teorema de Bjerknes resolve.
  \item \textbf{Giro da brisa}: no HS o vetor gira anti-horário ao
        longo do dia (espelho do HN).
  \item \textbf{Lenticular como nuvem que anda}: está PARADA sobre a
        crista com vento forte através — ótimo choque conceitual.
  \item \textbf{Föhn como ``compressão''}: o ar não é comprimido como
        pistão; desce pelo $\Gamma_d$ porque a chuva removeu água e
        $L_v$ a barlavento (T5 fecha a conta).
  \item \textbf{Vendaval de sotavento como ``vento que despenca''}: é
        transição hidráulica (Fr), não queda livre — a analogia da
        pia é a ponte certa.
\end{itemize}

\subsection*{Demonstração de baixo custo}

A pia da cozinha: jato supercrítico + salto hidráulico em anel —
todos já viram, ninguém tinha nomeado. Tanque com água e obstáculo
(régua) puxado a velocidades diferentes: os três regimes de Froude
aparecem. Vídeos: lenticulares em timelapse; Bora em Trieste.

\subsection*{Problemas sugeridos do livro (exercícios do Cap.~17)}

Aquecimento: $u_{eq}$ catabático com números dados; $\lambda_z$ de
ondas. Aprofundamento: föhn com LCL dado; Froude de camadas dadas.
Síntese: ``desenhe o dia de ventos de uma cidade costeira com serra
atrás (Santos!): brisa, terral, catabático e onda no mesmo ciclo''.
\emph{Confira a numeração na sua cópia (v1.02b).}

\subsection*{Uso do notebook \texttt{cap17\_ventosregionais.ipynb}}

\begin{center}
\begin{tabular}{lll}
\hline
Caso & Conteúdo & Momento sugerido \\
\hline
1 & brisa: EDO com Coriolis + hodógrafo diurno & Aula 1 \\
2 & catabático: EDO e $u_{eq}(\alpha, \Delta\theta)$ & Aula 1 \\
3 & ondas de montanha lineares (morro-sino) & Aula 2 \\
4 & água rasa sobre a serra: regimes de Froude & Aula 2 \\
\hline
\end{tabular}
\end{center}

\subsection*{Exercícios complementares e soluções}

Nas notas: T1--T5 e N1--N4. Soluções em \texttt{cap17\_solucoes.ipynb}
(\emph{não distribuir}): T2 com \texttt{sympy}; N1 responde quando a
brisa perde do vento sinótico; N4 constrói o mapa de regimes. Pares:
T2$\leftrightarrow$N2, T3$\leftrightarrow$N3, T4$\leftrightarrow$N4.

\subsection*{Conexões com o resto do curso}

Circuitos térmicos e Bowen (Cap.~3); $N$ como mola (Cap.~5);
lenticulares e rotores (Cap.~6); brisa como gatilho de convecção
(Cap.~14); catabático e piscina fria noturna (Cap.~18); arrasto de
onda nos modelos (Cap.~20); brisas e ilha de calor urbana (Cap.~19).
